\section{Scores by dataset}
\label{app:dataset-scores}

The figures show every dataset and all eleven methods on the original metric scale. Each score
marker is the mean of 250 outer-fold scores. For MCB, we average the five model seeds and five
outer folds within each chemical partition, giving 10 paired complete-partition scores. For
residual mean square (MSE), the one-step Hsu decision threshold is
$H=q_{1,11}\sqrt{\mathrm{MSE}/10}$, where $q_{1,11}$ is the corresponding 95\% MCB critical
value for eleven methods with 90 residual degrees of freedom. The pale region starts at the empirical winner and
extends by $H$ toward worse scores. A method inside the region cannot be ruled out as best; a method
beyond its boundary is classified as significantly worse by Hsu's test.

\begin{figure}[!htbp]
  \centering
  \includegraphics[width=\linewidth,height=.86\textheight,keepaspectratio]{native_score_hsu_decision_bars_paired_partition_average_primary_all.pdf}
  \caption{Full-grid TDC leaderboard metric values for all 22 ADMET datasets. The arrow below each
  heading gives the metric direction. Filled circles mark the empirical best, open circles mark other
  methods not ruled out as best, and crosses mark methods classified as significantly worse. The
  pale region shows the MCB decision threshold described above.}
  \label{fig:tdc-dataset-scores}
\end{figure}

\begin{figure}[!htbp]
  \centering
  \includegraphics[width=\linewidth,height=.86\textheight,keepaspectratio]{native_score_hsu_decision_bars_paired_partition_average_brier.pdf}
  \caption{Full-grid temperature-calibrated Brier scores for all 13 binary ADMET datasets. Lower is
  better. Marks and the decision region are defined as in Figure~\ref{fig:tdc-dataset-scores}. A common
  scalar temperature is fitted on each dedicated calibration holdout.}
  \label{fig:binary-dataset-scores}
\end{figure}

\clearpage
\begin{figure}[!htbp]
  \centering
  \includegraphics[width=\linewidth,height=.86\textheight,keepaspectratio]{native_score_hsu_decision_bars_paired_partition_average_interval_score.pdf}
  \caption{Full-grid method-specifically calibrated 90\% interval scores for all nine regression
  ADMET datasets. Lower is better. Marks and the decision region are defined as in
  Figure~\ref{fig:tdc-dataset-scores}. The horizontal axes differ because interval score is reported
  in the original units of each endpoint.}
  \label{fig:regression-dataset-scores}
\end{figure}

\FloatBarrier

Under the assumed block model, Hsu's selection rule includes at least one method with the best
population mean score with probability at least 95\% \citep[p.~1141]{hsu1984constrained}.
Excluding any particular method is a 5\%-level test of whether it performs worse than the best
alternative. If several methods share the optimal population mean, the 95\% guarantee concerns
including a best method, not simultaneously including all of them.

\subsection{Additional metrics on their original scales}

The following figures complement the mean-rank rails with the underlying full-grid metric values.
They use the same fixed method order and pairing as the three formal score sheets. Outlined
markers identify empirical bests. Horizontal axes are fitted separately for each dataset because
several regression metrics have
different endpoint units and scales.

\begin{figure*}[!htbp]
  \centering
  \includegraphics[width=\textwidth,height=.76\textheight,keepaspectratio]{appendix_native_scores_binary_auroc.pdf}
  \caption{Temperature-calibrated AUROC on all 13 binary datasets. Higher is better. Each marker
  is the mean of 250 outer-fold scores. Outlined markers identify empirical bests.}
  \label{fig:appendix-native-binary-auroc}
\end{figure*}

\begin{figure*}[p]
  \centering
  \includegraphics[width=\textwidth,height=.88\textheight,keepaspectratio]{appendix_native_scores_binary_average_precision.pdf}
  \caption{Temperature-calibrated average precision on all 13 binary datasets. Higher is better.
  Each marker is the mean of 250 outer-fold scores. Outlined markers identify empirical bests.}
  \label{fig:appendix-native-binary-ap}
\end{figure*}

\begin{figure*}[p]
  \centering
  \includegraphics[width=\textwidth,height=.88\textheight,keepaspectratio]{appendix_native_scores_binary_nll.pdf}
  \caption{Temperature-calibrated binary negative log-likelihood on all 13 binary datasets.
  Lower is better. Each marker is the mean of 250 outer-fold scores. Outlined markers identify
  empirical bests.}
  \label{fig:appendix-native-binary-nll}
\end{figure*}

\begin{figure*}[p]
  \centering
  \includegraphics[width=\textwidth,height=.88\textheight,keepaspectratio]{appendix_native_scores_binary_ece.pdf}
  \caption{Temperature-calibrated 10-bin expected calibration error on all 13 binary datasets.
  Lower is better. Each marker is the mean of 250 outer-fold scores. Outlined markers identify
  empirical bests.}
  \label{fig:appendix-native-binary-ece}
\end{figure*}

\begin{figure*}[p]
  \centering
  \includegraphics[width=\textwidth,height=.88\textheight,keepaspectratio]{appendix_native_scores_regression_mae.pdf}
  \caption{Mean absolute error on all nine regression datasets under the fixed method-specific
  calibration pipelines. Lower is better. Each marker is the mean of 250 outer-fold scores. The
  outlined markers identify empirical bests.}
  \label{fig:appendix-native-regression-mae}
\end{figure*}

\begin{figure*}[p]
  \centering
  \includegraphics[width=\textwidth,height=.88\textheight,keepaspectratio]{appendix_native_scores_regression_rmse.pdf}
  \caption{Root mean squared error on all nine regression datasets under the fixed
  method-specific calibration pipelines. Lower is better. Each marker is the mean of 250
  outer-fold scores. Outlined markers identify empirical bests.}
  \label{fig:appendix-native-regression-rmse}
\end{figure*}

\begin{figure*}[p]
  \centering
  \includegraphics[width=\textwidth,height=.88\textheight,keepaspectratio]{appendix_native_scores_regression_spearman.pdf}
  \caption{Spearman correlation on all nine regression datasets under the fixed method-specific
  calibration pipelines. Higher is better. Each marker is the mean of 250 outer-fold scores. The
  outlined markers identify empirical bests.}
  \label{fig:appendix-native-regression-spearman}
\end{figure*}

\begin{figure*}[p]
  \centering
  \includegraphics[width=\textwidth,height=.88\textheight,keepaspectratio]{appendix_native_scores_regression_nll.pdf}
  \caption{Gaussian negative log-likelihood on all nine regression datasets under the fixed
  method-specific calibration map. Lower is better. Each marker is the mean of 250 outer-fold
  scores. Outlined markers identify empirical bests.}
  \label{fig:appendix-native-regression-nll}
\end{figure*}

\begin{figure*}[p]
  \centering
  \includegraphics[width=\textwidth,height=.88\textheight,keepaspectratio]{appendix_native_scores_regression_coverage_error.pdf}
  \caption{Absolute 90\% coverage error, $|\mathrm{PICP}-0.90|$, where PICP is the prediction
  interval coverage probability, on all nine regression datasets under the fixed method-specific
  calibration map. Lower is better. Each marker is the
  mean of 250 outer-fold scores. Outlined markers identify empirical bests.}
  \label{fig:appendix-native-regression-coverage-error}
\end{figure*}
\FloatBarrier
